\subsection{System overview}\label{sec:overview}

The pipeline processes the video of a remotely operated vehicle (ROV)
between its
camera and the pilot. It was built for a BlueROV2 \cite{bluerov2}, whose onboard
Raspberry Pi runs BlueOS and the ArduSub autopilot \cite{ardusub} and
is joined to a
topside computer by an Ethernet tether, and it runs on any GNU/Linux machine. It
consists of two programs (Figure~\ref{fig:overview}). \eeye{}, the
vision layer, takes
frames from a camera or from the tether, applies a configurable chain of image
operators and an optional continuous visibility correction, and
publishes the result
(\emph{pilot}) and the untouched camera frame (\emph{raw}) on a shared-memory
\emph{frame bus} for any other program on the machine; it has no network code.
\hellion{} does everything that crosses the tether: it encodes the
video as H.264
\cite{itu_h264} over RTP \cite{rfc3550,rfc6184}, receives it topside,
relays it to
QGroundControl \cite{qgroundcontrol}, records it, serves a browser
control page, and
carries \eeye{}'s settings from the page to wherever \eeye{} runs.
The two programs
communicate only through the frame bus and \eeye{}'s configuration
file, so either can
be restarted without the other.

\eeye{} runs on one side of the tether, chosen at installation (the
  \emph{placement};
Figure~\ref{fig:overview}). In placement~B it runs on the vehicle
next to the camera,
the vehicle sends the processed picture, and the topside relays the
received packets to
QGroundControl unchanged, so the picture is encoded once. In
placement~A the vehicle
sends its camera, \eeye{} runs topside on the decoded stream, and the
processed picture
is encoded a second time for QGroundControl; this moves the image
processing off the
vehicle's processor. Autonomy code is expected to run next to \eeye{} and read
\emph{raw} from the bus, so that its input is the camera's frame and carries the
camera's timestamps.

\begin{figure*}[!t]
  \centering
  \begin{tikzpicture}[node distance=4mm and 5mm]
    % Placement B
    \node[lab, anchor=west] (bt) at (0,0) {\textbf{Placement B}};
    \node[ext, below=3mm of bt.south west, anchor=north west] (cam) {camera};
    \node[proc, right=of cam] (ee) {\eeye{}\\effects};
    \node[buf, right=of ee] (bus) {frame bus\\\emph{pilot}, \emph{raw}};
    \node[proc, right=of bus] (hr) {\hellion{}\\encode};
    \node[proc, below=14mm of hr] (ht) {\hellion{}\\relay, decode};
    \node[ext, right=10mm of ht, yshift=5mm] (qgc) {QGroundControl\\UDP 5600};
    \node[ext, right=10mm of ht, yshift=-6mm] (rec) {MP4 recording};
    \node[buf, left=of ht] (tbus) {frame bus\\\emph{pilot}};
    \node[ext, left=of tbus] (page) {browser page,\\webcam device};
    \draw[flow] (cam) -- (ee);
    \draw[flow] (ee) -- (bus);
    \draw[flow] (bus) -- (hr);
    \draw[net] (hr) -- node[lab, left, pos=0.45] {H.264/RTP\\over the
    tether} (ht);
    \draw[flow] (ht.east) -- ++(4mm,0) |- (qgc.west);
    \draw[flow] (ht.east) -- ++(4mm,0) |- (rec.west);
    \draw[flow] (ht) -- (tbus);
    \draw[flow] (tbus) -- (page);
    \begin{scope}[on background layer]
      \node[box, fill=blue!3, fit=(cam)(ee)(bus)(hr),
      label={[lab]above right:ROV}] {};
      \node[box, fill=green!3, fit=(page)(tbus)(ht)(qgc)(rec),
      label={[lab]below right:topside}] {};
    \end{scope}
    % Placement A
    \node[lab, anchor=west] (at) at (0,-5.4) {\textbf{Placement A}};
    \node[ext, below=3mm of at.south west, anchor=north west] (camA) {camera};
    \node[proc, right=of camA] (hrA) {\hellion{}\\encode};
    \node[proc, below=14mm of hrA] (htA) {\hellion{}\\decode};
    \node[buf, right=of htA] (busA) {frame bus\\\emph{camera}};
    \node[proc, right=of busA] (eeA) {\eeye{}\\effects};
    \node[buf, right=of eeA] (busA2) {frame bus\\\emph{pilot}};
    \node[proc, below=6mm of busA2] (heA) {\hellion{}\\encode};
    \node[ext, left=of heA] (qgcA) {QGroundControl, page,\\webcam device};
    \draw[flow] (camA) -- (hrA);
    \draw[net] (hrA) -- node[lab, left, pos=0.45] {H.264/RTP\\over
    the tether} (htA);
    \draw[flow] (htA) -- (busA);
    \draw[flow] (busA) -- (eeA);
    \draw[flow] (eeA) -- (busA2);
    \draw[flow] (busA2) -- (heA);
    \draw[flow] (heA) -- (qgcA);
    \begin{scope}[on background layer]
      \node[box, fill=blue!3, fit=(camA)(hrA), label={[lab]above right:ROV}] {};
      \node[box, fill=green!3, fit=(htA)(busA)(eeA)(busA2)(heA)(qgcA),
      label={[lab]below right:topside}] {};
    \end{scope}
  \end{tikzpicture}
  \caption{The two programs and the two placements. In placement B, used on the
    vehicle, \eeye{} processes the camera on the ROV, and
    QGroundControl receives the
    ROV's packets unchanged. In placement A the ROV sends its camera and the
    processing runs topside, at the cost of a second encode. Programs
    on one machine
  exchange frames only through the frame bus.}\label{fig:overview}
\end{figure*}
